\documentclass[12pt,a4paper]{abntex2}

% Pacotes de codificação e linguagem
\usepackage[utf8]{inputenc}
\usepackage[T1]{fontenc}
\usepackage[brazil]{babel}
\usepackage{csquotes}

% Pacotes de formatação
\usepackage{geometry}
\geometry{a4paper, margin=3cm, top=3cm, bottom=2cm}
\usepackage{setspace}
\usepackage{indentfirst}
\usepackage{times}
\usepackage{graphicx}
\usepackage{float}
\usepackage{abntex2cite}

% Espaçamento entre linhas
\onehalfspacing

% Dados do trabalho
\titulo{Relatório Técnico: Arquitetura de Soluções Industriais, Integração e Dados na WEG}
\autor{Lucas Miguel Leite}
\local{Sorocaba}
\data{06/02/2026}
\instituicao{
  ESCOLA E FACULDADE DE TECNOLOGIA SENAI GASPAR RICARDO JÚNIOR \\
  DESENVOLVIMENTO DE SISTEMAS
}
\tipotrabalho{Relatório técnico}
\preambulo{Relatório técnico apresentado à disciplina de Integração Vertical e Horizontal do Curso de Análise e Desenvolvimento de Sistemas, como requisito parcial para avaliação.}

% Professor
\newcommand{\professor}{Professor: Deivison Takatu}

% Palavras-chave
\palavraschave{WEG, Indústria 4.0, Integração de Sistemas, IoT Industrial, Manufatura Digital.}

\begin{document}

% Capa
\imprimircapa

% Folha de rosto
\imprimirfolhaderosto

% Resumo
\begin{resumao}
  O presente relatório técnico tem como objetivo analisar a arquitetura de soluções industriais da WEG, com foco nos fundamentos tecnológicos, na integração de sistemas e na gestão de dados no contexto da Indústria 4.0. A metodologia empregada baseia-se na revisão técnica das plataformas de hardware e software da companhia, investigando a convergência entre as tecnologias de operação (OT) e de informação (IT). O estudo descreve os mecanismos de sensoriamento e IoT, exemplificados pelo uso do sensor WEG Motor Scan e drives inteligentes, bem como a aplicação de Edge Computing para o processamento de dados na borda. São abordadas as estratégias de integração de sistemas corporativos através de softwares MES (Manufacturing Execution Systems), provenientes da aquisição da PPI-Multitask, que permitem a conexão vertical entre o chão de fábrica e sistemas ERP. Além disso, examina-se o tratamento de dados industriais massivos via plataforma WEGnology e a aplicação de Inteligência Artificial para manutenção preditiva com as soluções da BirminD. Conclui-se que a WEG consolidou uma transição estratégica de fornecedora de componentes para provedora de ecossistemas digitais integrados, oferecendo soluções escaláveis que aumentam a eficiência operacional e a confiabilidade dos ativos industriais.

  \vspace{1cm}
  \textbf{Palavras-chave:} WEG. Indústria 4.0. Integração de Sistemas. IoT Industrial. Manufatura Digital.
\end{resumao}

% Sumário
\pdfbookmark[0]{\contentsname}{toc}
\tableofcontents
\newpage

% Conteúdo Principal
\section{INTRODUÇÃO}

A WEG, multinacional brasileira consolidada na fabricação de equipamentos eletroeletrônicos, passou por uma transformação significativa na última década, migrando de uma fornecedora de hardware puro (motores e tintas) para uma provedora de soluções digitais completas para a Indústria 4.0.

O objetivo deste relatório é analisar os fundamentos tecnológicos que compõem o ecossistema digital da WEG, detalhando como a empresa realiza a integração de sistemas, gerencia dados industriais e conecta o chão de fábrica aos sistemas corporativos. A análise baseia-se na arquitetura de suas soluções recentes, como o WEG Motor Scan, a plataforma WEGnology e os sistemas de execução de manufatura (MES).

\section{FUNDAMENTOS TECNOLÓGICOS}

A base tecnológica da WEG estrutura-se na capacidade de coletar dados diretamente da fonte de energia e movimento.

\subsection{Sensoriamento e IoT (Internet das Coisas)}

O principal componente de entrada de dados no ecossistema atual é o WEG Motor Scan. Trata-se de um sensor retrofit (acoplável a motores antigos ou novos) que monitora vibração, temperatura e tempo de funcionamento.

\begin{itemize}
  \item \textbf{Tecnologia:} Utiliza comunicação Bluetooth\textsuperscript{\textregistered} e Gateway para envio de dados à nuvem.
  \item \textbf{Aplicação:} Permite a manutenção preditiva, substituindo a preventiva baseada apenas em horas de uso. O sensor capta o ``espectro de vibração'', fundamental para identificar desbalanceamentos e desalinhamentos mecânicos antes da falha \cite{weg2024}.
\end{itemize}

\subsection{Acionamentos e Controle}

No nível de automação, os Inversores de Frequência (série CFW) e Soft-Starters (série SSW) atuam não apenas como reguladores de potência, mas como sensores ativos.

\begin{itemize}
  \item \textbf{Funcionalidade:} Os inversores modernos da WEG possuem PLCs (Controladores Lógicos Programáveis) incorporados ou cartões de expansão que permitem processamento local de lógica, funcionando como a camada de ``automação distribuída''.
\end{itemize}

\section{INTEGRAÇÃO DE SISTEMAS}

A integração na arquitetura WEG segue o modelo de pirâmide de automação, mas com forte tendência ao ``achatamento'' via IoT.

\subsection{Protocolos e Interoperabilidade}

Para garantir a integração entre os equipamentos de chão de fábrica (OT) e os sistemas de gestão, a WEG utiliza padrões abertos de comunicação industrial.

\begin{itemize}
  \item \textbf{Redes de Campo:} Profibus DP, DeviceNet e CANopen são suportados nativamente nos drives e PLCs (Série PLC300/500).
  \item \textbf{Ethernet Industrial:} A migração para Profinet, Ethernet/IP e Modbus-TCP é predominante, permitindo maior largura de banda para tráfego de dados massivos.
  \item \textbf{OPC UA:} Para integração moderna, o padrão OPC UA é utilizado para comunicação segura e independente de plataforma entre o hardware de controle e sistemas SCADA ou MES.
\end{itemize}

\subsection{Edge Computing}

A WEG implementa soluções de Edge Computing (Computação de Borda) através de Gateways IoT (como o WEG Motor Scan Gateway). Isso permite que dados críticos sejam processados localmente para latência mínima, enquanto apenas dados agregados ou alarmes são enviados para a nuvem, otimizando o consumo de banda e armazenamento.

\section{SISTEMAS CORPORATIVOS E MES}

Um ponto de inflexão na estratégia da WEG foi a aquisição da PPI-Multitask e, mais recentemente, o investimento em empresas de software, consolidando a divisão WEG Digital Solutions.

\subsection{Sistemas MES (Manufacturing Execution Systems)}

Através das soluções advindas da PPI-Multitask, a WEG oferece sistemas MES que realizam o apontamento da produção em tempo real.

\begin{itemize}
  \item \textbf{Função:} O sistema coleta dados dos PLCs e sensores (OEE -- Overall Equipment Effectiveness) e elimina o apontamento manual em papel.
  \item \textbf{Integração Vertical:} O MES atua como a camada intermediária (``middleware industrial''), traduzindo os sinais elétricos das máquinas em ordens de produção compreensíveis para o ERP.
\end{itemize}

\subsection{Integração IT/OT}

A integração com sistemas corporativos (ERP) como SAP, TOTVS ou Oracle é realizada via APIs ou tabelas de interface. Isso permite que, quando uma ordem de venda é gerada no ERP, ela seja desdobrada automaticamente em ordens de produção no MES da WEG, que então configura os parâmetros das máquinas no chão de fábrica.

\section{DADOS INDUSTRIAIS E ANALYTICS}

A gestão de dados industriais na WEG culmina na plataforma WEGnology e no WEG Motion Fleet Management.

\subsection{WEG Motion Fleet Management (MFM)}

É a solução em nuvem dedicada ao gerenciamento da frota de ativos industriais.

\begin{itemize}
  \item \textbf{Visualização:} Dashboards apresentam o status de saúde dos motores (verde, amarelo, vermelho).
  \item \textbf{Arquitetura de Dados:} Os dados brutos (vibração nos eixos X, Y, Z) são enviados para a nuvem da WEG, onde algoritmos processam essas informações para gerar insights.
\end{itemize}

\subsection{Inteligência Artificial Industrial (BirminD)}

Com a aquisição da BirminD, a WEG integrou Inteligência Artificial (Machine Learning) aos seus dados industriais.

\begin{itemize}
  \item \textbf{Análise Preditiva Avançada:} A tecnologia permite não apenas ver o estado atual, mas prever a ``Vida Útil Remanescente'' (RUL) de um ativo.
  \item \textbf{Otimização:} Algoritmos analisam correlações complexas em bancos de dados industriais (Big Data) para sugerir otimizações de setpoint que economizam energia e reduzem refugo.
\end{itemize}

\section{CONCLUSÃO}

A análise dos fundamentos tecnológicos e da integração de sistemas da WEG demonstra uma estratégia madura de Indústria 4.0. A empresa superou a barreira de hardware integrando verticalmente a cadeia de valor: do sensor (Motor Scan) e do atuador (Inversor), passando pelo controle (PLC) e pela camada de borda (Gateway), até o software de gestão (MES/PPI-Multitask) e análise de dados em nuvem (WEGnology/BirminD).

Conclui-se que a arquitetura da WEG favorece a modularidade e a escalabilidade, permitindo que indústrias de diversos portes digitalizem suas operações sem a necessidade de substituição total do parque instalado, utilizando sensores retrofit e plataformas agnósticas de integração.

% Referências Bibliográficas
\bibliography{referencias}

\end{document}
